\section{Описание}

В работе реализован поиск одного образца в тексте при помощи Z-функции. В отличие от классического
символьного поиска, элементами последовательности являются не отдельные символы, а целые слова.
Поэтому Z-функция строится над массивом \texttt{std::vector<std::string>}. Это соответствует заданному
алфавиту: сравнение выполняется между словами, а не между буквами внутри слов.

Так как алфавит регистронезависимый, каждое считанное слово приводится к нижнему регистру. После этого
образец и текст объединяются в одну последовательность:
\[
P_1, P_2, \ldots, P_m,\ \$,\ T_1, T_2, \ldots, T_n.
\]
Разделитель \texttt{\$} не является словом латинского алфавита, поэтому он не может совпасть ни с одним
словом образца или текста. Для объединённой последовательности вычисляется Z-функция. Если для позиции
текстовой части значение Z-функции равно длине образца, то с этой позиции начинается очередное
вхождение.

\subsection{Приведение регистра}

Функция \texttt{toLowerCase} проходит по символам слова и применяет \texttt{std::tolower}. В функцию
передаётся \texttt{unsigned char}, чтобы избежать неопределённого поведения на отрицательных значениях
\texttt{char}, а результат явно приводится к \texttt{char} через \texttt{static\_cast}. Для латинских
слов этого достаточно, поскольку по условию требуется регистронезависимое сравнение.

\subsection{Построение Z-блоков}

Для позиции \(i\) значение \(Z[i]\) равно длине наибольшего префикса всей последовательности, который
совпадает с подпоследовательностью, начинающейся в \(i\). Алгоритм поддерживает текущий Z-блок
\([j, k)\) --- самый правый найденный отрезок, совпадающий с префиксом, где \(k\) является первым
индексом \emph{за} блоком (исключительная правая граница).

Если позиция \(i\) лежит внутри текущего блока (\(i < k\)), начальное значение берётся из уже
посчитанного симметричного значения:
\[
Z[i] = \min(k - i,\ Z[i - j]).
\]
Затем значение расширяется прямыми сравнениями слов, пока они совпадают. Если новый блок заканчивается
правее текущего (\(Z[i] > k - i\)), границы обновляются: \(j = i\), \(k = i + Z[i]\). Каждое слово
попадает в расширение правой границы ограниченное число раз, поэтому время построения Z-функции
линейно от длины объединённой последовательности.

\subsection{Чтение текста и позиции}

При чтении текстовой части программа сохраняет два массива. В \texttt{text} лежат нормализованные слова,
а в \texttt{info} для каждого слова хранится его исходная позиция: номер строки и номер слова в этой
строке. Пустые строки корректно учитываются как строки текста, но не добавляют слов.

После вычисления Z-функции обход начинается с первой позиции текстовой части, то есть с индекса
\texttt{patterns.size() + 1}. Если найдено совпадение длины \(m\), начало вхождения находится в массиве
позиций по индексу:
\begin{center}
\texttt{i - patterns.size() - 1}
\end{center}
Эта позиция и выводится в требуемом формате \texttt{строка, слово}.

Асимптотика реализации:
\begin{itemize}
\item чтение и нормализация входа -- \(O(m + n)\), где \(m\) -- число слов в образце, \(n\) -- число слов в тексте;
\item построение объединённой последовательности -- \(O(m + n)\);
\item вычисление Z-функции -- \(O(m + n)\);
\item память -- \(O(m + n)\) на хранение слов, позиций и массива Z.
\end{itemize}

\section{Исходный код}

Ниже приведены основные фрагменты реализации.

\lstinputlisting[
language=C++,
firstline=7,
lastline=12,
caption={Нормализация регистра слова},
label={lst:lower}
]{../main.cpp}

\lstinputlisting[
language=C++,
firstline=13,
lastline=33,
caption={Вычисление Z-функции над массивом слов},
label={lst:z-function}
]{../main.cpp}

\lstinputlisting[
language=C++,
firstline=36,
lastline=63,
caption={Чтение образца, текста и исходных позиций слов},
label={lst:input}
]{../main.cpp}

\lstinputlisting[
language=C++,
firstline=64,
lastline=78,
caption={Объединение образца с текстом и вывод найденных вхождений},
label={lst:search}
]{../main.cpp}
\pagebreak

\section{Консоль}
\begin{alltt}
$ g++ -std=c++17 -O2 -Wall -Wextra -pedantic main.cpp -o lab4.out
$ ./lab4.out
cat dog cat dog bird
CAT dog CaT Dog Cat DOG bird CAT
dog cat dog bird
1, 3
1, 8
\end{alltt}

В примере образец состоит из пяти слов. Благодаря приведению к нижнему регистру вхождения находятся
независимо от регистра слов во входном тексте.
\pagebreak
